\section{Formulación del problema}


¿Cómo desarrollar una metodología basada en métodos de ciencia de datos  para pronosticar la frecuencia y el tiempo de duración de interrupciones eléctricas (SAIDI y SAIFI) a nivel de zonas en redes de distribución de una empresa de energía eléctria, ubicada en Norte de Santander, identificando los factores que más influyen en su comportamiento.?

A partir de la formulación anterior, se derivan las siguientes preguntas de sistematización:

\begin{enumerate}
    \item ¿Cómo consolidar un repositorio analítico unificado y trazable que integre SCADA, OMS, GIS, CIS, registros de mantenimiento y datos climáticos, acompañado de un catálogo de datos y métricas de calidad que permitan para pronosticar la frecuencia y el tiempo de duración de interrupciones eléctricas?
    
    \item ¿Qué patrones temporales, espaciales, estacionales o comportamientos atípicos caracterizan la variabilidad de SAIDI y SAIFI a nivel de circuito, y de qué manera estos patrones influyen en la comprensión del fenómeno de continuidad del servicio elétrico y en la identificación de los factores que explican sus fluctuaciones?
    
    \item ¿Qué modelos tienen mejor desempeño para pronosticar la frecuencia y el tiempo de duración de interrupciones eléctricas (SAIDI y SAIFI) a nivel de circuito o zona?
    
    \item ¿Cómo construir artefactos reproducibles (notebooks, contenedores, documentación) y reportes técnicos alineados con los entregables del proyecto que faciliten la interpretación operativa y el uso de resultados en priorización de mantenimiento y reportes regulatorios?
\end{enumerate}


